\chapter{Planeamento das substituições por MDP e algoritmo genético}
\label{chap:substituicoes}

O módulo de planeamento de substituições de PLANTIA adapta a ordem dos blocos da rotação do módulo de planeamento de rotações
através de um MDP de horizonte finito, resolvido por programação dinâmica
retrospetiva, e compara essa solução exata com uma política adaptativa
obtida por algoritmo genético (AG). A especificação segue
\path{docs/SUBSTITUTION_PLANNING.md} e
\path{docs/GENETIC_SUBSTITUTIONS.md}. O MDP define o problema sequencial;
a programação dinâmica e o AG são dois métodos para procurar uma política
nesse mesmo problema.
O planeador reorganiza alinhamentos existentes dentro de cada período,
preservando os minutos individuais atribuídos nesse período.

\section{Estados e ações admissíveis}

O estado contém o índice do bloco temporal, os identificadores dos blocos
ainda disponíveis no período, o número de blocos consecutivos de cada
jogador em campo, o bloco anteriormente utilizado e um contexto observado.
O contexto assume um de três valores: equilibrado, criação de jogo ou
ressalto. Esta representação permite determinar ações, recompensa e
transições sem guardar todo o histórico do encontro.

Uma ação escolhe um dos blocos originais ainda disponíveis no período.
O bloco original zero é obrigatório no início, preservando o cinco inicial.
No período seguinte entram os blocos originalmente atribuídos a esse
período. Ações com o mesmo conjunto de jogadores usam um representante.
São rejeitadas escolhas que ultrapassem o limite de blocos consecutivos
ou que impossibilitem completar a rotação.

A reorganização conserva posições, disponibilidade, dimensão da rotação
e minutos por jogador e período. As pausas entre períodos não reiniciam
os contadores de utilização consecutiva. A implementação exata suporta
blocos de quatro, seis ou doze minutos; blocos mais finos são rejeitados
para limitar o crescimento da procura por permutações.

\section{Transições e recompensa}

A ocupação do campo, os blocos restantes e os contadores evoluem
deterministicamente com a ação. O contexto seguinte segue, por omissão,
a matriz
\[
P=\begin{pmatrix}
0{,}6&0{,}2&0{,}2\\
0{,}2&0{,}6&0{,}2\\
0{,}2&0{,}2&0{,}6
\end{pmatrix}.
\]
Todos os episódios começam em contexto equilibrado. A matriz é independente
da ação e constitui uma hipótese de simulação, não uma transição aprendida
com dados de adversários. A utilização real exigiria uma regra validada
para observar o contexto durante o jogo.

O bónus \(B_i(c)\) é zero no contexto equilibrado, \(0{,}4\widehat{A}_i\)
na criação de jogo e \(0{,}4\widehat{R}_i\) no ressalto. Para o conjunto
ativo \(L(a)\), contagem anterior \(k_i\) e número de jogadores que entram
\(n_{\mathrm{ent}}\), a recompensa é
\[
r(s,a)=\frac{\Delta}{48}\sum_{i\in L(a)}
[U_i+B_i(c)]\max(0,1-fk_i)-\kappa n_{\mathrm{ent}},
\]
com \(f=0{,}08\) e \(\kappa=0{,}15\) por omissão. O cinco inicial não paga
custo de substituição. Um bloco no banco reinicia a fadiga aproximada;
esta regra não constitui um modelo fisiológico. O termo dos coeficientes
de titulares é omitido porque o cinco inicial permanece fixo.

\section{Solução exata por programação dinâmica}

A utilidade total não é descontada. Para o estado terminal \(s_T\),
\[
\begin{aligned}
V(s_T)&=0,\\
V(s)&=\max_{a\in\mathcal{A}(s)}\biggl[r(s,a)\\
&\quad+\sum_{c'}P(c'\mid c)V(s'(s,a,c'))\biggr].
\end{aligned}
\]
A implementação memoriza os valores e exporta uma ação maximizadora para
cada estado de decisão viável avaliado. A optimalidade aplica-se apenas
às ações restritas, à recompensa e à matriz de transição especificadas.
Não são representados diferencial de pontos, faltas, novas lesões,
prolongamentos ou mudanças do alinhamento adversário.

\section{Algoritmo genético para políticas adaptativas}

\subsection{Representação e admissibilidade}

Um cromossoma contém um gene inteiro por estado de decisão viável.
Para os estados \(s_1,\ldots,s_n\), representa-se o indivíduo por
\[
\boldsymbol{g}=(g_1,\ldots,g_n),\qquad
 g_j\in\{0,\ldots,|\mathcal{A}(s_j)|-1\}.
\]
O gene seleciona uma ação da lista admissível do respetivo estado,
definindo uma política determinista \(\pi_{\boldsymbol{g}}\). Como o estado
inclui o contexto observado, o indivíduo representa uma política adaptativa,
não um horário fixo. Assim, a comparação com a programação dinâmica
mantém a mesma capacidade de reagir aos contextos modelados.

Uma recursão estrutural elimina ações que não permitem terminar a rotação
respeitando o descanso. Esta verificação não consulta recompensas nem
valores ótimos de Bellman. As ações preservam os titulares originais,
as posições e os minutos individuais por período. A evolução utiliza
apenas genes legais e não recebe ações ou valores ótimos do MDP.
A documentação inclui um teste que verifica que a evolução não chama
o otimizador por programação dinâmica.

\subsection{População e operadores evolutivos}

A população inicial contém 40 cromossomas com genes legais escolhidos
aleatoriamente. Após a avaliação inicial, executam-se 40 gerações com
os seguintes operadores:
\begin{itemize}
\item Seleção de progenitores por torneios de três indivíduos.
\item Cruzamento uniforme: cada gene é escolhido de um dos progenitores
 com probabilidade de 0,5 para cada um.
\item Mutação independente por gene com probabilidade de 0,1, sorteando
 novamente uma ação legal. O sorteio pode conservar o valor anterior.
\item Elitismo: os dois melhores indivíduos transitam sem alterações.
\end{itemize}
Como cada posição do cromossoma mantém o seu domínio de ações admissíveis,
os operadores não introduzem ações ilegais. São realizadas cinco execuções
independentes, com sementes de 42 a 46, para observar a variabilidade
entre pesquisas.

\subsection{Função de aptidão e referência exata}

A aptidão é a utilidade total esperada exata da política sob a matriz de
transição assumida. A avaliação de um cromossoma calcula
\[
\begin{aligned}
V^{\pi}(s_T)&=0,\\
V^{\pi}(s)&=r(s,\pi(s))\\
&\quad+\sum_{c'}P(c'\mid c)V^{\pi}(s'(s,\pi(s),c')).
\end{aligned}
\]
Ao contrário da equação de Bellman, esta expressão não maximiza sobre
ações: avalia a ação fixada pelo gene. A aptidão é
\(F(\boldsymbol{g})=V^{\pi_{\boldsymbol{g}}}(s_0)\), com contexto inicial
equilibrado. A avaliação memoriza valores e reutiliza a aptidão de
cromossomas repetidos. A melhor execução é selecionada pela utilidade
esperada exata, antes das comparações por simulação.

O valor ótimo \(V^*(s_0)\) da programação dinâmica fornece um limite
superior para o mesmo modelo. A diferença
\[
\Delta_{\mathrm{AG}}=V^*(s_0)-F(\boldsymbol{g}_{\mathrm{melhor}})
\]
deve ser não negativa, admitindo tolerância numérica. O AG é estocástico
e não fornece, em geral, prova de optimalidade. A igualdade com a referência
exata permite reconhecer uma solução ótima apenas para aquela execução
e aquele problema finito.

\section{Protocolo de comparação e limitações}

A comparação alargada avalia quatro políticas: solução exata por programação
dinâmica, AG, ordem fixa do módulo de planeamento de rotações e política gulosa. Esta última
maximiza a recompensa imediata entre ações com conclusão viável.
Todas utilizam os mesmos estados, ações, transições e recompensas.
O experimento genético recalcula a referência exata com parâmetros iguais,
sem alterar os relatórios originais do experimento MDP.

Além das esperanças exatas, as quatro políticas são avaliadas nas mesmas
500 trajetórias simuladas de contexto, com semente 10042, separada das
sementes de evolução. As comparações emparelhadas e os respetivos intervalos
medem apenas erro de amostragem da simulação. Não incluem incerteza dos
modelos preditivos, dos parâmetros ou dos resultados desportivos reais.
Todas as rotações simuladas recebem verificações independentes de minutos,
posições, disponibilidade, titulares e descanso, incluindo os totais
exatos por período.

O experimento original do MDP usa semente 42 na simulação e estuda
persistência entre 0,3 e 0,8 e fadiga entre 0 e 0,16. A comparação genética
inclui um relatório adicional com penalização de fadiga nula. Este cenário
estuda sensibilidade a uma hipótese; não afirma que jogadores reais não
sofrem fadiga. Os resultados dos dois cenários genéticos são apresentados
no Capítulo~\ref{chap:experiencias}.

No cenário documentado existem 46 estados de decisão viáveis, mas apenas
três oferecem várias ações, originando oito cromossomas possíveis. Todas
as cinco execuções avaliam os oito. A convergência pode, por isso, ser
plana desde a geração inicial. A igualdade com a solução exata neste
espaço reduzido não demonstra maior eficiência ou escalabilidade do AG.

A enumeração estrutural dos estados, a aptidão exata e a representação
completa por estado limitam ambos os métodos. O AG mantém a restrição a
blocos de quatro, seis ou doze minutos. Espaços maiores exigiriam uma
representação compacta da política ou avaliação aproximada da aptidão.
Os tempos registados separam enumeração, evolução e resolução exata,
excluindo simulação, figuras e exportação; não constituem um estudo
controlado de desempenho computacional.

Os testes documentados cobrem reprodutibilidade, elitismo, legalidade dos
genes, independência dos valores ótimos, avaliação da aptidão e respeito
pelo limite superior exato. Para a programação dinâmica, incluem também
comparação com enumeração exaustiva num exemplo determinista pequeno.

\section{Consulta e integração}

O comando \path{scripts/plan_substitutions.py} produz política, métricas,
simulações e exemplo de rotação. O comando
\path{scripts/query_substitution.py} consulta uma decisão a partir do
estado observado e devolve jogadores que entram e saem, papéis e utilidade
restante. Estados ausentes produzem erro. Mudanças de plantel,
disponibilidade ou minutos exigem gerar novamente a rotação e a política.

O comando \path{scripts/compare_genetic_substitutions.py} executa a
comparação genética e guarda os resultados em
\path{reports/genetic_substitutions/}. Os artefactos incluem
\texttt{metrics.json}, resultados por semente em \texttt{runs.csv},
curvas de convergência, \texttt{policy.json}, simulações e uma trajetória
exemplificativa. O marcador \texttt{\_SUCCESS} confirma a conclusão das
verificações. A opção \texttt{--fatigue 0}, com outro diretório de saída,
reproduz a análise sem penalização de fadiga.

A política genética usa o mesmo formato de consulta:
\begin{quote}
\small
\path{python scripts/query_substitution.py}\\
\texttt{--policy}\\
\path{reports/genetic_substitutions/policy.json}
\end{quote}
A opção \texttt{--state} permite fornecer um estado observado posterior.
O campo \texttt{value} representa a utilidade esperada restante da política
selecionada, não uma afirmação geral de optimalidade.

É possível comparar políticas para um jogo do módulo de planeamento de vários jogos fornecendo o
respetivo diretório através de \texttt{--rotation-dir}. Esta execução
preserva os minutos selecionados e não substitui automaticamente a política
MDP do horizonte integrado. Uma eventual escolha da política genética
para utilização deve ser distinguida da sua avaliação comparativa.
